\scanpage{177}\mspage{32}%
% C original page filename: H134-05_134-2_scan177.tex
% archive sheet number: 177
où ``pris à la lettre'' $p''$ est maintenant une inclusion, mais ``moralement'' c'est \inserted{tjrs} une fibration à fibres très \unsure{polies|jolies} [de dim finie notamment, voire des variétés non singulières - c'est donc beaucoup mieux que ce que nous donne le yoga de Cartan-Serre ``toute application continue est équivalente à une fibration''\ldots] C'est ce dernier diagramme qui est susceptible d'une construction directe, purement algébrique, dans le cadre des schémas, lorsque l'on a développé la construction du voisinage tubulaire étale.

Le point où je veux en venir, c'est que la considération \inserted{de la cat} des faisceaux (ou \inserted{complexes d'iceux, ou} $n$-champs \unsure{..|\ldots}) $\unsure{\mathcal P|P}$-constructibles sur un $X$, où $\mathcal P$ est une stratification ``équisingulière'' donnée\raisebox{0pt}[0pt][0pt]{\parbox{0pt}{\figure{\begin{tikzpicture}[overlay,x=1mm,y=1mm]\fill (-6,-1.5) circle[radius=.12mm];\end{tikzpicture}}}}, revient dans le cas particulier actuel à la connaissance d'un diagramme de \uline{types d'homotopie}\raisebox{0pt}[0pt][0pt]{\parbox{0pt}{\figure{\begin{tikzpicture}[overlay,x=1mm,y=1mm]\fill (6,.8) ellipse[x radius=.08mm,y radius=.19mm];\end{tikzpicture}}}} \inserted{(ou de pro-types, si on revient à la top. étale)}\raisebox{0pt}[0pt][0pt]{\parbox{0pt}{\figure{\begin{tikzpicture}[overlay,x=1mm,y=1mm]\draw[line width=.2pt] (2.4,1.15) -- (2.5,.7);\end{tikzpicture}}}} ordinaires
% The parenthesis is physically above ordinaires, without an attachment guide.
% Its visible order before ordinaires is retained; ordinaires is not marked as an insertion.
% Nontext marks above: (1158,1368) below donnée; (1633,1538) right of underline;
% (1333,1563) just above/right of the parenthesis. Presence is clear; origin is uncertain.
\diagram{\begin{tikzcd}[ampersand replacement=\&,column sep=large,row sep=large]
\unsure{\Delta|\triangle} \arrow[r,"\unsure{J|j}"] \arrow[d,"p"'] \& X_1\\
X_0 \&
\end{tikzcd}}
\raisebox{0pt}[0pt][0pt]{\parbox{0pt}{\figure{\begin{tikzpicture}[overlay,x=1mm,y=1mm]\draw[line width=.3pt] (8,4.4) -- (8.05,4.2) -- (8,4);\end{tikzpicture}}}}en prenant des systèmes de coeff locaux (ou de $n$-champs loc. constants) sur ces sommets \inserted{$X_0,X_1$ \deleted{$\Delta$}} reliés par des hom de \unsure{compatibilité|compatibilit} du type $p^*(F_0)\to J^*(F_1)$. Ce serait un exercice amusant (que je n'\unsure{ai} pas encore fait) de
% The isolated short dark mark at (315,1849), above prenant, has uncertain origin.
% The crossed triangular sign is in the insertion above the gap after sommets;
% the tall live l of reliés passes into the same crowded area.
% The manuscript stops at the complete word de; no following text is supplied.
